\section{Supplementary material}

\subsection{Extended probe results}
\label{sec:suppl-probe}
\cref{tab:suppl-probe} gives the per-population and feature-count-matched probes for the arena (seed 0), the MLP
probe on all wirings, and the NYC dark-duration numbers. ``all@762'' is a random 762-neuron subset of the student,
the GRU's width; ``rand@CX'' is a random subset the size of the central complex.

\begin{table}[h]
  \centering\small
  \setlength{\tabcolsep}{3pt}
  \resizebox{\linewidth}{!}{%
\begin{tabular}{lcccccc}
    \toprule
    wiring & all & all@762 & CX & rand@CX & EPG & MLP, all \\
    \midrule
    connectome & 8.5 & 18.9 & 12.6 & 10.7 & 50.1 & 4.7\pm0.1 \\
    connectome, random signs & 7.9 & 16.8 & 11.7 & 10.4 & 44.0 & 4.7\pm0.3 \\
    block-shuffled & 11.1 & 20.9 & 14.4 & 13.2 & 38.5 & 5.9\pm0.3 \\
    degree-shuffled & 12.6 & 20.1 & 14.1 & 15.0 & 35.2 & 6.9\pm0.7 \\
    random & 12.8 & 19.2 & 14.5 & 14.4 & 37.3 & 6.4\pm0.4 \\
    connectome, unsigned & 12.8 & 20.6 & 13.9 & -- & 34.2 & -- \\
    dense & 21.5 & 28.8 & 25.2 & 23.7 & 39.8 & 41\pm24 \\
    GRU, parameter-matched & 31.6 & -- & -- & -- & -- & 27.4\pm0.3 \\
    GRU, recurrent-budget-matched & 35.2 & -- & -- & -- & -- & -- \\
    \bottomrule
  \end{tabular}}
  \caption{Median heading error (\degree) in the 8:16 blind window, arena, seed 0 (MLP column: mean\pm sd over 3
  seeds). The EPG population alone carries almost no heading; the central complex decodes no better than a random
  subset of its size.}
  \label{tab:suppl-probe}
\end{table}

\begin{table}[h]
  \centering\small
  \resizebox{\linewidth}{!}{%
\begin{tabular}{lcc}
    \toprule
    wiring (NYC) & 16 dark (8:24) & 20 dark (4:24) \\
    \midrule
    connectome & 12.2\pm2.2, 0.88 & 9.5\pm1.3, 0.88 \\
    block-shuffled & 16.4\pm1.5, 0.86 & 12.2\pm2.0, 0.84 \\
    degree-shuffled & 25.2\pm1.8, 0.69 & 15.9\pm0.7, 0.73 \\
    random & 21.6\pm1.1, 0.80 & 14.0\pm0.3, 0.83 \\
    GRU, parameter-matched & 49.4\pm1.9, 0.08 & 39.8\pm1.5, 0.25 \\
    \bottomrule
  \end{tabular}}
  \caption{Dark duration on the NYC twin: heading error (\degree) and position $R^2$, ridge probe, 3 seeds.}
  \label{tab:suppl-dark-nyc}
\end{table}

\subsection{Goal-vector navigation}
\label{sec:suppl-goalnav}
\begin{table}[h]
  \centering\small
  \setlength{\tabcolsep}{3pt}
  \resizebox{\linewidth}{!}{%
\begin{tabular}{llcccc}
    \toprule
    wiring & env & readout & no cut & cut 3 & cut 5 \\
    \midrule
    connectome & arena & oracle goal vector & 0.92 / 0.70 & 0.92 / 0.70 & 0.92 / 0.70 \\
    connectome & arena & random actions & 0.14 / 0.12 & 0.14 / 0.12 & 0.14 / 0.12 \\
    connectome & arena & state (ridge) & 0.22 / 0.16 & 0.22 / 0.15 & 0.22 / 0.14 \\
    connectome & arena & state (MLP) & 0.24 / 0.15 & -- & -- \\
    block-shuffled & arena & state & 0.24 / 0.15 & 0.23 / 0.14 & 0.21 / 0.13 \\
    degree-shuffled & arena & state & 0.22 / 0.13 & 0.23 / 0.14 & 0.21 / 0.13 \\
    random & arena & state & 0.21 / 0.13 & 0.21 / 0.12 & 0.19 / 0.12 \\
    dense & arena & state & 0.20 / 0.12 & 0.21 / 0.12 & 0.22 / 0.12 \\
    GRU, param.-matched & arena & state & 0.22 / 0.15 & 0.22 / 0.14 & 0.22 / 0.16 \\
    GRU, budget-matched & arena & state & 0.23 / 0.14 & 0.23 / 0.11 & 0.23 / 0.13 \\
    \midrule
    connectome & NYC, greedy & oracle & 0.28 / 0.25 & -- & 0.28 / 0.25 \\
    connectome & NYC, greedy & state & 0.05 / 0.05 & -- & 0.04 / 0.04 \\
    \midrule
    connectome & NYC, bumper & oracle & 0.83 / 0.78 & 0.83 / 0.78 & 0.83 / 0.78 \\
    connectome & NYC, bumper & random actions & 0.00 / 0.00 & 0.00 / 0.00 & 0.00 / 0.00 \\
    connectome & NYC, bumper & state (ridge) & 0.18 / 0.18 & 0.21 / 0.21 & 0.17 / 0.17 \\
    connectome & NYC, bumper & state (MLP) & 0.27 / 0.25 & -- & -- \\
    block-shuffled & NYC, bumper & state & 0.10 / 0.10 & 0.14 / 0.14 & 0.13 / 0.13 \\
    random & NYC, bumper & state & 0.06 / 0.05 & 0.02 / 0.02 & 0.04 / 0.04 \\
    transformer & NYC, bumper & state & 0.16 / 0.15 & 0.14 / 0.14 & 0.16 / 0.15 \\
    GRU, param.-matched & NYC, bumper & state & 0.11 / 0.11 & 0.11 / 0.11 & 0.12 / 0.12 \\
    GRU, param.-matched & NYC, bumper & state (MLP) & 0.17 / 0.16 & -- & -- \\
    \midrule
    connectome & NYC dusk / overcast / fog / night & state & 0.23 / 0.20 / 0.17 / 0.13 & -- & -- \\
    block-shuffled & NYC dusk / overcast / fog / night & state & 0.13 / 0.10 / 0.10 / 0.00 & -- & -- \\
    transformer & NYC dusk / overcast / fog / night & state & 0.13 / 0.17 / 0.13 / 0.17 & -- & -- \\
    GRU, param.-matched & NYC dusk / overcast / fog / night & state & 0.13 / 0.13 / 0.20 / 0.13 & -- & -- \\
    \bottomrule
  \end{tabular}}
  \caption{Goal-vector navigation: success / SPL~\cite{anderson2018evaluation}. Arena: 50 episodes per seed, 3 seeds.
  NYC: 30 episodes per seed, 3 seeds for state readouts, goals 30 to 60\,m. The vision cut lasts 3 or 5 steps from
  step 6. With the greedy controller even the oracle goal vector fails on streets (it walks into buildings), so the NYC
  rows use a bumper controller (on collision: turn two steps toward the goal side, commit three steps forward), which
  lifts the oracle to 0.83. Under it the state readouts clear random actions and random wiring, the connectome is the
  best core (0.18 ridge, 0.27 MLP, seed 0), and the vision cut barely matters: the readout, not vision loss, is the
  limit. The last block is the clear-trained readout evaluated under appearance shifts (seed 0, success only).}
  \label{tab:goalnav}
\end{table}

\subsection{Data efficiency}
\label{sec:suppl-data}
\begin{table}[h]
  \centering\small
  \resizebox{\linewidth}{!}{%
\begin{tabular}{lcccc}
    \toprule
    train episodes & connectome & block-shuffled & degree-shuffled & GRU \\
    \midrule
    50 & 0.970\pm0.009 & 0.989\pm0.009 & 0.955\pm0.013 & \textbf{0.821\pm0.007} \\
    100 & 0.899\pm0.003 & 0.924\pm0.006 & 0.913\pm0.003 & \textbf{0.762\pm0.016} \\
    200 & 0.796\pm0.011 & 0.818\pm0.021 & 0.791\pm0.009 & \textbf{0.655\pm0.007} \\
    500 & 0.677\pm0.003 & 0.702\pm0.015 & 0.713\pm0.004 & \textbf{0.573\pm0.002} \\
    1000 & 0.579\pm0.003 & 0.585\pm0.005 & 0.610\pm0.002 & \textbf{0.434\pm0.004} \\
    \bottomrule
  \end{tabular}}
  \caption{Blind MSE (arena, mean over horizons, 3 seeds; 5 at 1000) against training-set size. PCA basis fit on the
  full training set in every row.}
  \label{tab:data}
\end{table}

\subsection{Subgraph ablation}
\begin{table}[h]
  \centering\small
  \resizebox{\linewidth}{!}{%
\begin{tabular}{lccc}
    \toprule
    subgraph & connectome & degree-shuffled & margin \\
    \midrule
    full graph ($\ge$5 synapses, 91,663 edges) & 0.575 & 0.610 & 0.035 \\
    $\ge$1 synapse & 0.618 & 0.661 & 0.043 \\
    $\ge$10 synapses & 0.584 & 0.604 & 0.020 \\
    central complex only & 0.588 & 0.608 & 0.020 \\
    mushroom bodies removed & 0.582 & 0.609 & 0.027 \\
    descending neurons removed & 0.588 & 0.610 & 0.022 \\
    \bottomrule
  \end{tabular}}
  \caption{Blind MSE 8:16, arena, seed 0, real wiring against the degree-shuffled rewiring of the same subgraph. The
  full-graph margin is 0.035\pm0.005 over 5 seeds.}
  \label{tab:ablation}\label{sec:ablation}
\end{table}

\subsection{Real handheld video}
\label{sec:suppl-realvid}
The clips are walked with a phone held at chest height; each is centre-cropped to a square, resampled to 2\,Hz and cut
into non-overlapping 24-step episodes. Yaw per step is the mean horizontal phase-correlation shift over a $4\times4$
patch grid converted through a 60\degree\ horizontal field of view; forward motion is the mean radial expansion. Steps
with $|{\rm yaw}| > 6^\circ$ become turn actions, steps with expansion above 1\,px forward, below $-1$\,px back, the
rest stay. The heading target of the probe is the integrated continuous yaw, not the quantised record, so the
``actions only'' floor and the probe measure different things on this data. Zero-shot, fine-tuned and probe numbers for
every architecture are in \cref{tab:realvid}; per-clip numbers are in \texttt{results/summary.md}.

\subsection{Teacher-free planning (negative)}
\label{sec:suppl-plan}
Latent CEM-MPC with the student as dynamics model succeeds in 0 to 3\,\% of arena navigation episodes for every
wiring, while the same planner with oracle dynamics succeeds in 100\,\% and a scripted policy in 100\,\%; on the NYC
twin student-model planning is at 0\,\% under clear, dusk, night, fog and overcast appearance. The students' 8-step
rollouts are not accurate enough to plan with. Dropped from the claims.

\subsection{Video}
The supplementary video (\texttt{wm2cx\_paper.mp4}, generated by \texttt{scripts/make\_video\_paper.py}) shows one
arena episode side by side for the \vjepa teacher, the connectome-wired student and the parameter-matched GRU: what each
receives (panorama, fly-eye input, neuron or hidden-unit activity), what each thinks (its latent guess on the teacher's
2-D latent trajectory, with the cosine match over time) and where each believes it is (the pose decoded from its own
state, the view rendered there, and the decoded path on the top-down map). The released strip figures are cut from it. A second
video (\texttt{wm2cx\_realvid.mp4}, \texttt{scripts/make\_video\_realvid.py}) repeats the comparison on a held-out real
phone clip with both students fine-tuned on the other clips.

\subsection{Reproducibility}
All results in the paper come from one run on a single RTX 4090 (about 14 GPU-hours): 5 seeds on the arena, 3 on NYC,
12 wirings, 424 checkpoints. Code, environments, the ingest pipeline for phone video, and the checkpoints are released;
the \flywire-derived connectome file is rebuilt by one script from the public v783 release.
